The digital transformation of retail banking has shifted the volume and
velocity of financial transactions well beyond what manual oversight can
monitor. In a network of 150 branches processing thousands of teller
operations daily, the gap between transaction speed and supervisory
capacity creates exposure to fraud, money laundering, operational errors,
and behavioral anomalies that may go undetected until periodic compliance
reviews surface them --- often weeks or months after the fact.

Amen Bank's Direction Centrale du Système d'Information already operates
an AI-based anti-money laundering system oriented toward bank-wide
compliance reporting. This internship addresses a complementary need:
real-time decision support for branch supervisors at the point of
transaction. The objective is to design and build a platform that ingests
teller operations as they occur, profiles both the client and the
employee involved, detects anomalies using deep learning models, and
surfaces explainable alerts that a non-technical supervisor can
understand and act upon immediately.

The platform covers four teller-window operations --- retrait, versement,
virement, and remise de chèques --- and monitors two independent actors
per transaction: the client initiating it and the employee processing it.
An Autoencoder detects point anomalies in individual transactions while
an LSTM network identifies sequential deviations in client behavior over
time. Every alert is accompanied by a human-readable explanation
generated through SHAP, bridging the gap between model output and
supervisory action.

Beyond the detection models themselves, the project invests in the
infrastructure that makes a machine learning system operational and
maintainable: a streaming pipeline built on RedPanda, behavioral
profiles maintained through a hot/cold data layer, automated experiment
tracking with MLflow, three independent retraining triggers, structured
logging, and Prometheus-based observability. The system runs as 20
Docker containers orchestrated through Docker Compose, validated
end-to-end on a synthetic dataset of 243,000 transactions generated
from five client archetypes calibrated against the bank's published
data.

This rapport is organized as follows. Chapter~1 presents the
institutional context, the regulatory framework governing anti-money
laundering in Tunisia, and a review of the anomaly detection techniques
that inform the system design. Chapter~2 describes the overall
architecture, the data schemas, and the key design decisions.
Chapter~3 covers the synthetic data generation pipeline, including
client archetype design and anomaly scenario injection. Chapter~4
details the training and evaluation of the Autoencoder and LSTM models.
Chapter~5 presents the streaming pipeline, from event ingestion through
enrichment, scoring, and decision-making, including batch integration
and the dead letter queue. Chapter~6 covers the MLOps layer: experiment
tracking, model promotion, retraining triggers, the supervisor feedback
loop, and observability. Chapter~7 describes the Streamlit dashboard
and presents the end-to-end validation results. Chapter~8 documents the
production enhancements --- Kubernetes deployment, CI/CD, secrets
management, load testing, and data retention --- that would be required
to transition the platform from its current development deployment to a
production banking environment.